% =====================================================================
%  TCET CODING PLATFORM -- Official Technical Documentation
%  MODULE 6 of 7 : Proctoring & Academic Integrity
%  STANDALONE FILE.  Compile: pdflatex Module6_Proctoring.tex  (twice)
% =====================================================================
\documentclass[12pt,a4paper,oneside]{report}
\usepackage[T1]{fontenc}
\usepackage[utf8]{inputenc}
\usepackage[a4paper,margin=1in]{geometry}
\usepackage{mathptmx}\usepackage{microtype}\usepackage{graphicx}\usepackage{xcolor}
\usepackage{tabularx}\usepackage{booktabs}\usepackage{longtable}\usepackage{xltabular}
\usepackage{array}\usepackage{enumitem}\usepackage{listings}\usepackage{titlesec}
\usepackage{fancyhdr}\usepackage{parskip}\usepackage[hidelinks]{hyperref}
\definecolor{tcetblue}{RGB}{0,0,0}\definecolor{tcetaccent}{RGB}{0,0,0}
\definecolor{codebg}{RGB}{245,246,248}\definecolor{codeframe}{RGB}{215,219,225}
\definecolor{codekw}{RGB}{0,0,0}\definecolor{codecomment}{RGB}{0,0,0}
\definecolor{codestr}{RGB}{0,0,0}\definecolor{boxframe}{RGB}{0,0,0}
\hypersetup{colorlinks=true,linkcolor=tcetblue,urlcolor=tcetaccent,citecolor=tcetblue,
  pdftitle={TCET Coding Platform -- Module 6: Proctoring},pdfauthor={Centre of Excellence, TCET}}
\lstdefinestyle{tcp}{backgroundcolor=\color{codebg},basicstyle=\ttfamily\footnotesize,
  keywordstyle=\color{codekw}\bfseries,commentstyle=\color{codecomment}\itshape,
  stringstyle=\color{codestr},numberstyle=\tiny\color{gray},frame=single,
  rulecolor=\color{codeframe},breaklines=true,showstringspaces=false,tabsize=2,
  columns=fullflexible,keepspaces=true,xleftmargin=0.6em,framexleftmargin=0.6em}
\lstset{style=tcp}
\newcommand{\code}[1]{\texttt{#1}}
\newenvironment{callout}[1]{\par\vspace{4pt}\noindent
  \begin{tabular}{@{}p{0.014\linewidth}@{\hspace{0.6em}}p{0.93\linewidth}@{}}
  \textcolor{boxframe}{\rule{2pt}{\baselineskip}} & \textbf{\textcolor{tcetblue}{#1}}\\[2pt] & }{\end{tabular}\par\vspace{6pt}}
\titleformat{\chapter}[display]{\bfseries\fontsize{18}{22}\selectfont}{\bfseries\fontsize{14}{17}\selectfont Chapter \thechapter}{0.2em}{}
\titlespacing*{\chapter}{0pt}{6pt}{22pt}
\titleformat{\section}{\bfseries\fontsize{16}{19}\selectfont}{\thesection}{0.6em}{}
\titleformat{\subsection}{\bfseries\fontsize{14}{17}\selectfont}{\thesubsection}{0.5em}{}
\pagestyle{fancy}\fancyhf{}\renewcommand{\headrulewidth}{0.4pt}
\fancyhead[L]{\small\itshape TCET Coding Platform}\fancyhead[R]{\small\itshape Module 6 \textbullet\ Proctoring}
\fancyfoot[C]{\thepage}
\fancypagestyle{plain}{\fancyhf{}\fancyfoot[C]{\thepage}\renewcommand{\headrulewidth}{0pt}}
\newcommand{\note}[1]{\par\vspace{3pt}\noindent{\color{tcetaccent}\textbf{Note.}}~#1\par\vspace{3pt}}
\begin{document}
\begin{titlepage}\centering
\vspace*{1.4cm}{\color{tcetaccent}\rule{\linewidth}{1.2pt}}\\[0.5cm]
{\LARGE\bfseries TCET Coding Platform\par}
\vspace{0.4cm}{\Large Official Technical Documentation\par}{\large Draft~1 \textemdash\ Testing Phase\par}
{\color{tcetaccent}\rule{\linewidth}{1.2pt}}\\[1.4cm]
{\Huge\bfseries\color{tcetblue} Module 6\par}\vspace{0.3cm}
{\Huge Proctoring \& Academic Integrity\par}\vspace{2.2cm}
{\large\bfseries Powered by\par}{\Large\color{tcetblue} Centre of Excellence (CoE), TCET\par}\vfill
\begin{tabular}{rl}\textbf{Module} & 6 of 7\\[3pt]
\textbf{Scope} & Client enforcement, server scoring, integrity guarantees\\[3pt]
\textbf{Status} & Draft~1 (living document until Stage~2)\\[3pt]
\textbf{Audience} & Officials, Security reviewers, Faculty\\\end{tabular}
\vspace{0.8cm}{\small\itshape \today\par}\end{titlepage}
\pagenumbering{roman}\tableofcontents\clearpage\pagenumbering{arabic}

% =====================================================================
\chapter{Overview \& Philosophy}
% =====================================================================
\section{Scope of this Module}
Live contests are protected by a layered proctoring and academic-integrity
subsystem that combines browser-level enforcement, server-side event logging, and
scoring penalties. This module specifies what is prevented, what is detected and
penalized, how enforcement is made authoritative on the server, and \textemdash\
in the interest of an honest official record \textemdash\ what browser-based
proctoring can and cannot guarantee.

\section{Threat Model}
The design accepts that a web page cannot fully control the operating system and
therefore adopts a two-tier philosophy.
\begin{description}[leftmargin=1.6em,style=nextline]
  \item[Prevent what the browser can prevent] Actions the browser can block
        outright (copy / cut / paste, right-click, text selection, dangerous
        shortcuts) are suppressed and cost the student nothing, because the
        action never succeeded.
  \item[Detect and penalize what it cannot] Actions the browser cannot hard-block
        (leaving the tab, leaving fullscreen, taking a screenshot) are detected,
        logged for faculty review, and scored as violations.
\end{description}

\begin{callout}{A deliberate distinction}
Only unpreventable evasions move the violation counter. Blocking a copy already
protects the content, so charging a penalty for it would punish a failed action.
Leaving the tab cannot be blocked, so it is detected and penalized. The scoring
model reflects this distinction precisely.
\end{callout}

% =====================================================================
\chapter{Client-Side Enforcement}
% =====================================================================
\section{Activation}
A frontend proctoring controller activates for the duration of an active attempt
and installs a coordinated set of controls, removing them cleanly when the
attempt ends or the student leaves the workspace.

\section{Fullscreen Lock}
The contest is intended to run in fullscreen. Leaving fullscreen is recorded as a
violation and immediately triggers an attempt to re-enter fullscreen. While the
page is not in fullscreen, a lock overlay covers the contest, and the very next
user interaction anywhere on the page is used as the gesture the browser requires
to restore fullscreen \textemdash\ so the student is returned to fullscreen
without having to find a button.

\section{Focus, Tab \& Visibility Detection}
Losing window focus or tab visibility (switching apps, minimizing, or moving to
another tab) is detected, the contest surface is \emph{blanked} instantly so a
capture taken while the browser is in the background contains nothing readable,
and the event is recorded as a violation.

\section{Screenshot Mitigation}
A screenshot key press is intercepted (on both key-down and key-up, since
platforms differ in which they deliver), recorded as a violation, and the system
clipboard is immediately overwritten \textemdash\ the only lever a web page has
over a screenshot that has already been taken.

\section{Blocked Shortcuts \& Actions}
The controller suppresses shortcuts that would open another surface, leave the
page, or reveal developer tools \textemdash\ refresh, fullscreen toggle,
developer tools, and new-tab / new-window / print / save / view-source
combinations. It also blocks the right-click context menu, copy / cut / paste,
and selection or dragging of question text (the code editor manages its own
selection and is exempt). A leave-page guard warns against accidental navigation
away from a live attempt.

\section{Debounced Reporting}
A single physical action often fires several browser events (one key can emit
both a fullscreen change and a focus loss). Events sharing a logical bucket
within a short cooldown are reported once, so one student action costs exactly one
violation \textemdash\ never several.

\section{Visual Deterrents}
The contest interface additionally displays a candidate watermark and a screen
guard overlay, discouraging photography of the screen and making any leaked
capture traceable.

\section{Summary of Client Controls}
\begin{longtable}{@{}p{0.30\linewidth} p{0.30\linewidth} p{0.30\linewidth}@{}}
\toprule
\textbf{Control} & \textbf{Mechanism} & \textbf{Effect}\\
\midrule
\endhead
Fullscreen lock & Fullscreen API + overlay & Re-enters fullscreen; penalizes exit.\\
Focus/visibility blanking & Focus \& visibility events & Blanks contest; penalizes leaving.\\
Screenshot mitigation & Key interception + clipboard wipe & Penalizes and neutralizes capture.\\
Shortcut suppression & Key handling & Prevents escape/devtools/print.\\
Clipboard \& menu block & Event prevention & Prevents extraction of content.\\
Watermark \& screen guard & Overlay & Deters and traces photography.\\
\bottomrule
\end{longtable}

% =====================================================================
\chapter{Server-Side Recording \& Enforcement}
% =====================================================================
\section{The Server is Authoritative}
Client-side controls are advisory; the server is authoritative. Every proctoring
event is sent to the backend, which persists an immutable record for faculty
review, determines whether the event is scored, applies the penalty, and, if the
maximum-violations threshold is reached, ends the attempt as auto-submitted
\textemdash\ auto-submitting any pending draft for judging, exactly like a manual
submit.

\begin{callout}{Why server-side matters}
Tampering with the client cannot award a student points or hide a completed
attempt: scoring, the violation count, and the auto-submit threshold are all
computed on the backend from persisted state. The client can only make the
experience smoother; it cannot change the outcome.
\end{callout}

\section{Scored vs. Logged Events}
\begin{longtable}{@{}p{0.28\linewidth} p{0.14\linewidth} p{0.48\linewidth}@{}}
\toprule
\textbf{Event} & \textbf{Scored?} & \textbf{Meaning}\\
\midrule
\endhead
Tab switch & Yes & Window lost focus / switched away.\\
Visibility loss & Yes & Document became hidden.\\
Fullscreen exit & Yes & Left fullscreen mode.\\
Screenshot & Yes & Screenshot key detected.\\
Copy / Cut / Paste & No & Clipboard action blocked (logged only).\\
Context menu & No & Right-click blocked (logged only).\\
\bottomrule
\end{longtable}

\section{The Event Log}
Proctoring events are stored immutably with the event type, the attempt and
contest, the student, a timestamp, and a human-readable detail. The log gives
faculty an auditable trail of what happened during any attempt, independent of
the scoring outcome.

% =====================================================================
\chapter{Integrity Guarantees}
% =====================================================================
\section{Guarantees During a Live Contest}
Taken together with the deferred-grading model of Module~5, the subsystem
provides the following guarantees during a live contest.
\begin{longtable}{@{}p{0.34\linewidth} p{0.58\linewidth}@{}}
\toprule
\textbf{Guarantee} & \textbf{Basis}\\
\midrule
\endhead
No answer leakage & Nothing is graded until faculty publishes.\\
No clipboard exfiltration & Copy / cut / paste and context menu are blocked.\\
Screen-capture resistance & Leaving the window blanks the contest; screenshots wipe the clipboard and are penalized.\\
Enforced focus & Leaving fullscreen or the tab is penalized and auto-corrected; repeated evasion ends the attempt.\\
Complete, fair grading & Even a disconnected student is finalized at their real deadline with their last work judged.\\
Auditable trail & Every event is recorded immutably for review.\\
\bottomrule
\end{longtable}

\section{Interaction with Scoring}
Violations reduce the attempt score by a fixed penalty each and, at the
threshold, end and grade the attempt. Because scoring is deferred and
server-side, a student never learns during the contest how a violation affected
their score, preserving the integrity of the examination.

% =====================================================================
\chapter{Known Limitations}
% =====================================================================
\section{What Browser Proctoring Cannot Do}
In the interest of an honest official record, the platform documents the inherent
limits of any purely browser-based proctor.
\begin{itemize}[leftmargin=1.5em]
  \item It cannot prevent an external camera from photographing the screen
        (mitigated, not prevented, by the watermark and screen guard).
  \item It cannot control a second physical device or a second computer.
  \item It cannot defeat operating-system screen capture that bypasses the
        browser entirely; it can detect the standard key path and blank on focus
        loss, but not all capture routes.
\end{itemize}

\section{Compensating Controls}
\begin{longtable}{@{}p{0.34\linewidth} p{0.58\linewidth}@{}}
\toprule
\textbf{Limitation} & \textbf{Compensating control}\\
\midrule
\endhead
External camera / photography & Candidate watermark and screen guard; physical invigilation.\\
Second device & Invigilated laboratory conduct; future integrations.\\
OS-level capture & Focus-loss blanking; penalty on the standard key path.\\
\bottomrule
\end{longtable}
These limitations motivate the future-scope items \textemdash\ invigilated /
lockdown environments and, where policy permits, webcam-assisted proctoring
\textemdash\ discussed in Module~7.

% =====================================================================
\chapter{Faculty Perspective \& Summary}
% =====================================================================
\section{What Faculty See}
Faculty can review, for any attempt, the full sequence of proctoring events, the
resulting violation count and penalty, and whether the attempt ended by manual
submit, deadline, or violation limit \textemdash\ alongside the graded outcome
and the student's final code.

\section{Recommended Conduct}
Because browser proctoring is strong but not absolute, the platform is designed
to be used within \emph{invigilated laboratory sessions}: the technical controls
raise the cost and detectability of evasion, while human invigilation covers the
physical channels the browser cannot. This combination is what makes the
examinations trustworthy at scale.

\section{Module Summary}
The proctoring subsystem prevents what the browser can prevent and detects and
penalizes what it cannot, records every event immutably, and enforces all scoring
consequences on the server. Combined with deferred grading, it gives fair,
auditable, tamper-resistant examinations, within clearly stated limits that
physical invigilation and future integrations are intended to close.

\vspace{0.6cm}
\begin{center}
\textit{--- End of Module 6. Continued in Module 7: Deployment, Infrastructure, Operations \& Future Scope. ---}
\end{center}
\end{document}
